
\medskip

%On doit appliquer deux couches de peinture sur le sol et les parois intérieures d'une piscine rectangulaire dont les dimensions sont données dans le document 2.

%À l'aide des documents ci-dessous, calculer le budget que l'on doit prévoir pour les travaux de peinture.
%
%\begin{center}
%\begin{tabularx}{\linewidth}{|m{5.6cm}|X|}\hline
%\textbf{Document 1 :} pot de peinture
%
%Surface pouvant être peinte: $35$ m$^2$
%
%Prix : \np{12000}~F&
%\textbf{Document 2 :} piscine de base rectangulaire 
%
%Longueur : 8~m Largeur : 4~m Profondeur : 1,70~m
%
%\begin{center}
%\psset{unit=0.8cm}
%\begin{pspicture}(0,-1.4)(7.8,3.5)
%%\psgrid
%\pspolygon(0,1.4)(5.8,0)(6.5,2.3)(0.7,3.7)
%\psline(0.7,3.7)(0.7,2.3)(6.1,1)
%\psline(0,1.4)(0,0)(5.8,-1.4)(5.8,0)
%\psline(5.8,-1.4)(6.5,0.9)(6.5,2.3)
%\psline(0.7,2.3)(0.4,1.3)
%\rput(3.6,3.4){8 m}
%\rput(0,2.8){4 m}
%\uput[r](6.5,1.6){1,70 m}
%\end{pspicture}
%\end{center}\\ \hline
%\end{tabularx}
%\end{center}
%
%\textbf{Toute trace de recherche même non aboutie sera prise en compte dans la notation.}
L'intérieur de la piscine est constitué de :

\begin{itemize}
\item deux rectangles de 8 m  sur 1,70 m ;
\item deux rectangles de 4 m sur 1,70 m ;
\item du fond de 8 m sur 4 m
\end{itemize}

L'aire de la surface à peindre est donc égale à :

$2 \times 8 \times 1,70 + 2 \times 4 \times 1,70 + 8 \times 4 = 27,2 + 13,6 + 32 = 72,8$~m$^2$.

Appliquer deux couches revient à peindre $2 \times 72,8 = 145,6$~m$^2$.

Or $\dfrac{145,6}{35} \approx 4,2$.

Quatre pots ne suffiront pas : il faut acheter 5 pots de peinture à \np{12000}~F soit un budget de \np{60000}~F{}.

\bigskip

